\begin{tikzpicture}[
    >=Latex,
    every node/.style={font=\small},
    layer/.style={
        rectangle, rounded corners=4pt, draw, very thick,
        minimum width=11.0cm, minimum height=1.4cm,
        align=center, inner sep=6pt,
    },
    composition/.style={layer, fill=purple!10, draw=purple!55!black},
    presentation/.style={layer, fill=blue!8,    draw=blue!55!black},
    interface/.style={layer, fill=red!8,        draw=red!55!black},
    domainlayer/.style={layer, fill=green!10,   draw=green!45!black},
    infrastructure/.style={layer, fill=gray!15, draw=gray!60!black},
    module/.style={
        rectangle, rounded corners=2pt, draw, thin,
        fill=white, minimum width=2.4cm, minimum height=0.65cm,
        align=center, font=\scriptsize\ttfamily, inner sep=2pt,
    },
    layerlabel/.style={
        font=\bfseries\small, anchor=west, inner sep=2pt,
    },
    dep/.style={->, very thick, blue!50!black},
    sidenote/.style={
        font=\scriptsize\itshape, gray!65!black,
        align=left, anchor=west,
    },
]
    % =========================================================
    % Schichten (von oben nach unten)
    % =========================================================
    \node[composition]                                  (L1) {};
    \node[presentation,    below=0.35cm of L1]          (L2) {};
    \node[interface,       below=0.35cm of L2]          (L3) {};
    \node[domainlayer,     below=0.35cm of L3]          (L4) {};
    \node[infrastructure,  below=0.35cm of L4]          (L5) {};

    % =========================================================
    % Schicht-Beschriftungen (links)
    % =========================================================
    \node[layerlabel, anchor=east]
        at ([xshift=-4pt]L1.west) {Komposition};
    \node[layerlabel, anchor=east]
        at ([xshift=-4pt]L2.west) {Präsentation};
    \node[layerlabel, anchor=east]
        at ([xshift=-4pt]L3.west) {Schnittstelle};
    \node[layerlabel, anchor=east]
        at ([xshift=-4pt]L4.west) {Domäne};
    \node[layerlabel, anchor=east]
        at ([xshift=-4pt]L5.west) {Infrastruktur};

    % =========================================================
    % Module pro Schicht
    % =========================================================
    % Komposition
    \node[module] at (L1) {app};

    % Präsentation: graphics, ui, controls
    \node[module] (m_graphics) at ([xshift=-3.0cm]L2.center) {graphics};
    \node[module] (m_ui)       at (L2.center)                {ui};
    \node[module] (m_controls) at ([xshift=3.0cm]L2.center)  {controls};

    % Schnittstelle: exporter
    \node[module] at (L3) {exporter};

    % Domäne: geometry, slicing
    \node[module] (m_geometry) at ([xshift=-1.6cm]L4.center) {geometry};
    \node[module] (m_slicing)  at ([xshift=1.6cm]L4.center)  {slicing};

    % Infrastruktur: core
    \node[module] at (L5) {core};

    % =========================================================
    % Abhängigkeitsrichtung (Pfeil rechts neben den Schichten)
    % =========================================================
    \draw[dep, very thick]
        ([xshift=0.6cm]L1.east) -- ([xshift=0.6cm]L5.east)
        node[midway, sloped, above, font=\scriptsize, gray!70!black]
        {Abhängigkeitsrichtung};

    % =========================================================
    % Erläuterungen rechts
    % =========================================================
    \node[sidenote] at ([xshift=1.3cm]L1.east)
        {kennt alle\\anderen Schichten};
    \node[sidenote] at ([xshift=1.3cm]L2.east)
        {Fenster, GUI,\\Eingabe};
    \node[sidenote] at ([xshift=1.3cm]L3.east)
        {JSON-Serialisierung};
    \node[sidenote] at ([xshift=1.3cm]L4.east)
        {Mesh, Slicing-\\Algorithmen};
    \node[sidenote] at ([xshift=1.3cm]L5.east)
        {Logging,\\Basisfunktionen};

\end{tikzpicture}